\input{estilo}
\usetikzlibrary{decorations.pathmorphing}
\begin{document}
\begin{figura}\begin{tikzpicture}[y=0.95cm,
    evento/.style={circle, fill=tinta, inner sep=0pt, minimum size=4.5pt},
    onda/.style={->, decorate, decoration={snake, amplitude=0.5mm, segment length=3.5mm, pre length=1mm, post length=2.5mm}}]
  \foreach \x/\p in {0/P, 4/Q, 8/R} {
    \draw (\x,0) -- (\x,7.6);
    \node[font=\small, above] at (\x,7.6) {proceso \var{\p}};
  }
  % eventos: nombre, x, y, lado del rótulo
  % eventos: proceso, número, x, y, lado del rótulo
  \foreach \l/\i/\x/\y/\lado in {p/1/0/0.7/left, p/2/0/2.7/left, p/3/0/4.8/left, p/4/0/6.8/left,
      q/1/4/0.75/right, q/2/4/1.8/right, q/3/4/2.75/right, q/4/4/3.8/below right,
      q/5/4/4.8/right, q/6/4/5.8/right, q/7/4/6.75/right,
      r/1/8/0.7/right, r/2/8/2.75/right, r/3/8/4.7/right, r/4/8/6.75/right} {
    \node[evento] (\l\i) at (\x,\y) {};
    \node[etiq, \lado=2pt] at (\l\i) {\var{\l}\textsubscript{\i}};
  }
  \draw[onda] (p1) -- (q2);
  \draw[onda] (q1) -- (p2);
  \draw[onda] (q1) -- (r4);
  \draw[onda] (q4) -- (r3);
  \draw[onda] (q5) -- (p4);
  \draw[onda] (r2) -- (q7);
  \draw[->, suave] (-1.2,0.3) -- node[etiq, sloped, above] {tiempo} (-1.2,2.3);
\end{tikzpicture}\end{figura}
\end{document}
